% ============================================================
%  SECTION 7: PROPOSED MODEL OF FUNCTIONAL MATURITY ASSESSMENT
% ============================================================

\section{Proposed Framework for Functional Maturity Assessment}

The preceding sections indicate that honey maturity is a continuous and multidimensional phenomenon. Therefore, no single parameter, whether morphological, microbiological, or physicochemical, is sufficient for its reliable assessment. Rather than treating honey as simply “mature” or “immature”, it is more appropriate to assess the degree of functional maturity in relation to microbiological stability, chemical quality, and fermentation risk. Natural variability related to botanical origin, climate, weather conditions, hive type, and beekeeping practices makes rigid universal classification difficult. However, the main quality and commercial risk associated with insufficient maturity, namely fermentation, can be assessed using measurable parameters, particularly water activity. The framework proposed in this section should be regarded as a conceptual and operational tool requiring further validation before possible use in official control practice.

% The considerations conducted in the preceding sections led to
% the formulation of the thesis that honey maturity is a continuous
% and multidimensional phenomenon, and that no single parameter --
% neither morphological (capping) nor microbiological (CFU/g) --
% is sufficient for its reliable assessment.
% Since honey maturity/ripening is a continuous phenomenon, only the
% concept of ``degree of maturity'' should be used. The introduction
% of any classification in this regard is not so much extremely
% difficult as practically impossible due to the natural variability
% of honey resulting from at least climatic, weather, and beekeeping
% practice conditions, and the diversity of nectar sources. However,
% the significant risk that honey maturity poses to product safety
% and commercial value -- namely the risk of fermentation -- can be
% limited. As indicated above, this risk is directly dependent on
% water activity.

% As indicated above, this risk is directly dependent on water activity. Accordingly, the functional maturity assessment model proposed in this section should be regarded as a conceptual and operational framework requiring further validation before possible use in official control practice.

% \subsection{Proposed biomarkers of functional maturity}

% The selection of biomarkers for the functional maturity assessment
% model is based on four criteria: (1) connection with a specific
% biological mechanism of honey ripening, (2) analytical validation
% or advanced validation stage, (3) technical and economic
% accessibility in control laboratories, (4) independence from
% climatic conditions and production system.

\subsection{Water Activity as a Priority Parameter}

Among all the biomarkers mentioned, water activity
($a_w$) is proposed as the priority parameter for three reasons.
First, it is a thermodynamic measure of the \emph{availability
of water to microorganisms}, not merely the water content per unit
mass of product; it directly correlates with the biological
possibility of osmophilic yeast growth, making it the only
parameter that can replace the entire range of indirect
microbiological indicators~\cite{ruegg1981,waaw2006}.
Second, measurement of $a_w$ is rapid (3--5 minutes), inexpensive,
and can be performed at an apiary or control laboratory using a
standard instrument such as AquaLab CX3TE or Rotronic HygroLab;
the result is reproducible and unambiguous, requiring no expert
interpretation~\cite{zamora2006}.
Third, the criterion $a_w < 0.60$ is based on a biologically
justified mechanism -- it is the physiological threshold of
osmophilic yeasts determined in vitro under isobaric
conditions~\cite{waaw2006,ruegg1981} -- and not on an arbitrarily
chosen administrative value.

It is nevertheless necessary to conduct studies on whether,
similarly to permitted exceptions regarding water content, exceptions
should also be established in this regard. The microbiological
stability of certain honey types may result not so much from low
water activity as from specific components present in plant nectar.

%,----------------------------------------------------------
\subsection{Hierarchical Decision Framework}
%,----------------------------------------------------------

Based on the biomarkers described above, the proposed functional maturity assessment framework can be structured into three analytical levels, organized according to the principle of proportionality. Each successive level is analytically more complex and should be applied only when the preceding level does not provide an unambiguous assessment or when the result is borderline.

\subsubsection{Level I -- Basic Screening}

Level I constitutes a first-order filter based on two parameters
measured directly on the product:

\begin{itemize}
    \item \textbf{Moisture $\leq 20\%$} taking into account
          legally specified exceptions (refractometry,
          Codex STAN 12-1981 method), a normative parameter,
          mandatory in all major standards;
    \item \textbf{Water activity $a_w < 0.60$} taking into
          account possible exceptions, a parameter proposed as
          a normative supplement; when both of these criteria
          are met, the honey is \emph{microbiologically stable
          by virtue of the state of microbiological knowledge},
          % #R1.5: "impossible" replaced — more cautious phrasing per reviewer request; supported by Ruegg & Blanc (1981) and WAAW (2006)
          and yeast growth is \sout{impossible} \textcolor{blue}{not observed and is effectively inhibited under standard storage conditions,}
          regardless of their count~\cite{ruegg1981,waaw2006}.
\end{itemize}

\textbf{Key rule of Level I:} if moisture $\leq 20\%$
\textsc{AND} $a_w < 0.60$, the product is not subject to
disqualification based on yeast count (CFU/g), regardless of the
determined value of this parameter. Disqualification based solely
on CFU/g, without confirmation of exceeding the $a_w$ threshold,
is scientifically unjustified and legally
inadequate~\cite{codex_stan12,eudir2001110}.

If Level I indicates a borderline state ($18.1$--$20.0\%$
moisture or $a_w = 0.60$--$0.62$) or the yeast count exceeds
$1000$~CFU/g, analysis is escalated to Level~II.

\subsubsection{Level II -- Extended Qualitative}

Level II provides an assessment of functional maturity in all
three dimensions defined in Section~3.3: microbiological, chemical,
and biological stability. Level~II parameters include:

\begin{enumerate}
    \item \textbf{Enzymatic profile:}
    \begin{itemize}
        \item Diastase activity $\geq 8$~Schade units (or $\geq 3$~units
              for botanical types with naturally low amylolytic
              activity, e.g., acacia, citrus, lavender)~\cite{codex_stan12};
    \end{itemize}
    \item \textbf{Chemical parameters:}
    \begin{itemize}
        \item F/G ratio $\geq 0.9$; HMF $\leq 40$~mg/kg;
              free acidity $\leq 50$~mEq/kg
              (Codex)~\cite{codex_stan12};
    \end{itemize}
\end{enumerate}

% #R2.6: NEW FIGURE — FMAM hierarchical decision flowchart added per reviewer request
\begin{figure}[ht]
\centering
\tikzset{
  fmbox/.style={rectangle, rounded corners=4pt, draw=black, fill=gray!10,
                text width=4.2cm, align=center, minimum height=0.7cm, font=\footnotesize},
  levelbox/.style={rectangle, rounded corners=4pt, draw=blue!70!black, fill=blue!8,
                   text width=5.0cm, align=center, minimum height=1.0cm, font=\footnotesize},
  okbox/.style={rectangle, rounded corners=4pt, draw=green!60!black, fill=green!8,
                text width=3.4cm, align=center, minimum height=0.7cm, font=\footnotesize},
  warnbox/.style={rectangle, rounded corners=4pt, draw=orange!80!black, fill=orange!8,
                  text width=3.4cm, align=center, minimum height=0.7cm, font=\footnotesize},
  fmarrow/.style={-{Stealth[length=2mm]}, thick}
}
\resizebox{\linewidth}{!}{%
\begin{tikzpicture}[node distance=0.6cm and 0.8cm]
  \node[fmbox] (start)
    {Honey sample};

  \node[levelbox, below=of start] (lI)
    {\textbf{Level I --- Basic Screening}\\
     Moisture $\leq 20\%$ \textsc{and} $a_w < 0.60$\\
     \scriptsize(refractometry + $a_w$ meter)};

  \node[okbox,  below left =1.2cm and 0.6cm of lI] (pass1)
    {Microbiologically stable\\No CFU/g disqualification};

  \node[warnbox, below right=1.2cm and 0.6cm of lI] (border)
    {Borderline or CFU/g\,$> 1000$\\$\Rightarrow$ escalate to Level~II};

  \node[levelbox, below=1.0cm of border] (lII)
    {\textbf{Level II --- Extended Qualitative}\\
     Diastase $\geq 8$\,SU;\; HMF $\leq 40$\,mg/kg\\
     F/G $\geq 0.9$;\; Acidity $\leq 50$\,mEq/kg};

  \node[okbox,  below left =1.2cm and 0.5cm of lII] (pass2)
    {Functionally mature\\Compliant};

  \node[warnbox, below right=1.2cm and 0.5cm of lII] (fail2)
    {Parameter(s) exceeded\\Further investigation};

  \draw[fmarrow] (start)  -- (lI);
  \draw[fmarrow] (lI) -- node[above left,  font=\scriptsize]{Both met}   (pass1);
  \draw[fmarrow] (lI) -- node[above right, font=\scriptsize]{Borderline / fail} (border);
  \draw[fmarrow] (border) -- (lII);
  \draw[fmarrow] (lII) -- node[above left,  font=\scriptsize]{All met}       (pass2);
  \draw[fmarrow] (lII) -- node[above right, font=\scriptsize]{Any exceeded}  (fail2);
\end{tikzpicture}}% end \resizebox
\caption{\textcolor{blue}{Hierarchical decision structure of the proposed Functional Maturity Assessment
Framework (FMAM). Level~I provides rapid basic screening based on moisture and water
activity; Level~II provides extended qualitative assessment for borderline cases. The framework
is a conceptual model requiring further multilaboratory validation prior to adoption in
official control practice.}}
\label{fig:fmam}
\end{figure}
% END NEW FIGURE

%,----------------------------------------------------------
\subsection{Recommendations for Regulatory Bodies and Apimondia}
%,----------------------------------------------------------

% #R1.5 #R2.6: "HDM" was inconsistent with the framework name used throughout; standardised to FMAM
Based on the proposed \sout{HDM} \textcolor{blue}{FMAM} and the conducted scientific analysis,
this article formulates the following recommendations addressed
to key stakeholders:

\subsubsection{Recommendations for Apimondia}

\begin{enumerate}
    \item \textbf{Supplement the 2023 position with the $a_w$
          criterion:} the Apimondia recommendation should
          explicitly state that the presence of osmophilic yeasts
          in honey does not constitute grounds for declaring
          immaturity, provided moisture $\leq 20\%$ and
          $a_w < 0.60$; both parameters should be listed as
          simultaneous necessary
          conditions~\cite{ruegg1981,waaw2006};
    \item \textbf{Take into account the regional diversity of
          production systems:} the position should contain a clear
          clause regarding honeys from tropical climates and
          traditional systems, recognizing technological
          dehumidification as a legal production practice~\cite{tadele2025}
          or excluding specific dehumidification techniques
          that negatively affect the product;
    \item \textbf{Initiate dialogue with ISO TC~34/SC~17:}
          the ISO 24607:2025 standard on bee honey should be
          extended to include the $a_w$ criterion as a normative
          parameter~\cite{codex_stan12}.
\end{enumerate}

\subsubsection{Recommendations for EU bodies (EC, DG SANTE)}


\textbf{Consider the inclusion of $a_w$  as a risk-based control parameter in official honey quality controls}, particularly for consignments with borderline moisture values, products from high-humidity production systems, or cases in which fermentation risk is analytically uncertain.


\subsubsection{Recommendations for laboratories and the scientific community}

\begin{enumerate}
    \item \textbf{Standardize the multiparametric microbiological
          assessment protocol} encompassing simultaneous measurement
          of CFU/g, $a_w$, temperature, and phase state of honey
          (liquid/crystallized); results should be reported
          jointly, not as isolated values~\cite{zamora2006,analytica2020}.
    \item \textbf{Perform multilaboratory validation of the
          decenedioic acid biomarker} and extend its verification
          beyond three botanical types (rapeseed, acacia, jujube)
          to tropical, stingless bee, and traditional production
          system honeys~\cite{sun2021}.
    \item \textbf{Develop portable in situ analysis tools:}
          miniaturized NIR and Raman spectrometers allow
          measurement of moisture, $a_w$, and sugar profile
          directly at the apiary or at the border, with accuracy
          comparable to laboratory
          methods~\cite{honeyid2025,schoder2026}.
\end{enumerate}
